%Tarea de clase 6
\documentclass[12pt]{beamer}
\setbeamercolor{title}{bg = magenta!50!black, fg= white}
\usepackage[utf8]{inputenc}
\usepackage[spanish]{babel}
\spanishdatedel
\uselanguage{Spanish} % Teoremas en idioma que elijas.
\languagepath{Spanish}
%% -----------------------------------------------------------------------------
%% 2017 por Fausto M. Lagos S. <piratax007@protonmail.ch>
%% 
%% Este trabajo puede ser distribuido o modificado bajo los
%% términos y condiciones de la LaTeX Project Public License (LPPL) v1.3C, 
%% o cualquier versión posterior. La última versión de esta licencia
%% puede verse en:
%% http://www.latex-project.org/lppl.txt
%% 
%% -----------------------------------------------------------------------------
%% Usted es libre de usarlo, modificarlo o distribuirlo siempre que se
%% respeten los términos de la licencia y se reconozca a los autores originales
%% -----------------------------------------------------------------------------
%% Estos comandos y ambientes están desarrollados para el Curso de LaTeX
%% de https://www.youtube.com/c/BrainOnTube y algunos han sido desarrollados
%% por https://github.com/jdleesmiller/latex-course, tenga presente sus
%% nombres y licencia para poder hacer uso de éstos.
%% -----------------------------------------------------------------------------

\usepackage{xkeyval}
\usepackage{minted}
\usepackage{tikz}
\usetikzlibrary{
	calc, angles, quotes, positioning,
    babel, shapes, arrows, shadows
    }
\usepackage{tcolorbox}
\usepackage{xstring}

\newcommand{\bftt}[1]{\textbf{\texttt{#1}}}
\newcommand*\keystroke[1]{%
  \tikz[baseline=(key.base)]
    \node[%
      draw,
      fill=white,
      drop shadow={shadow xshift=0.25ex,shadow yshift=-0.25ex,fill=black,opacity=0.75},
      rectangle,
      rounded corners=2pt,
      inner sep=1pt,
      line width=0.5pt,
      font=\scriptsize\sffamily
    ](key) {#1\strut}
  ;
}
\newcommand{\keystrokebftt}[1]{\keystroke{\bftt{#1}}}
\renewcommand{\comment}[1]{{\color[HTML]{008080}\textit{\textbf{\texttt{#1}}}}}
\newcommand{\cmd}[1]{{\color[HTML]{008000}\bftt{#1}}} % comandos
\newcommand{\bs}{\char`\\}
\newcommand{\cmdbs}[1]{\cmd{\bs#1}}
\newcommand{\lcb}{\char '173}
\newcommand{\rcb}{\char '175}
\newcommand{\cmdbegin}[1]{\cmdbs{begin\lcb}\bftt{#1}\cmd{\rcb}}
\newcommand{\cmdend}[1]{\cmdbs{end\lcb}\bftt{#1}\cmd{\rcb}}
\newcommand{\wllogo}{\textbf{Overleaf}}

\newenvironment{exampletwouptiny}
  {\VerbatimEnvironment
   \begin{VerbatimOut}{example.out}}
  {\end{VerbatimOut}
   \setlength{\parindent}{0pt}
   \fbox{\begin{tabular}{l|l}
   \begin{minipage}{0.55\linewidth}
     \inputminted[fontsize=\scriptsize,resetmargins]{latex}{example.out}
   \end{minipage} &
   \begin{minipage}{0.35\linewidth}
     \setlength{\parskip}{6pt plus 1pt minus 1pt}%
     \raggedright\scriptsize\input{example.out}
   \end{minipage}
   \end{tabular}}}
   
\newenvironment{exampletwouptinynoframe}
  {\VerbatimEnvironment
   \begin{VerbatimOut}{example.out}}
  {\end{VerbatimOut}
   \setlength{\parindent}{0pt}
   \begin{tabular}{l|l}
   \begin{minipage}{0.55\linewidth}
     \inputminted[fontsize=\scriptsize,resetmargins]{latex}{example.out}
   \end{minipage} &
   \begin{minipage}{0.35\linewidth}
     \setlength{\parskip}{6pt plus 1pt minus 1pt}%
     \raggedright\scriptsize\input{example.out}
   \end{minipage}
   \end{tabular}}

\newcommand\diag[4]{%
  \multicolumn{1}{p{#2}|}{\hskip-\tabcolsep
  $\vcenter{\begin{tikzpicture}[baseline=0,anchor=south west,inner sep=#1]
  \path[use as bounding box] (0,0) rectangle (#2+2\tabcolsep,\baselineskip);
  \node[minimum width={#2+2\tabcolsep-\pgflinewidth},
        minimum  height=\baselineskip+\extrarowheight-\pgflinewidth] (box) {};
  \draw[line cap=round] (box.north west) -- (box.south east);
  \node[anchor=south west] at (box.south west) {#3};
  \node[anchor=north east] at (box.north east) {#4};
 \end{tikzpicture}}$\hskip-\tabcolsep}}
 
\newtheorem{definition}{\bf Definición:}

\newtheorem{ejemplo}{\bf Ejemplo:}[section]

\newcounter{remark}
\newenvironment{remark}[1]
{
	\refstepcounter{remark}
	\begin{tcolorbox}[colback = myBlue!25, colframe = firebrick!75, title=\textbf{Observación \theremark: } #1, arc = 3mm, sharp corners = northwest]
	\fontfamily{qag}\selectfont
}
{
	\end{tcolorbox}
}

\newcommand{\Ref}[2]{
	\IfEqCase {#1}{
		{fig}{
			\textbf{Fig. \ref{#2}}
		}
		{tab}{
			\textbf{Tabla. \ref{#2}}
		}
	}
}
\usepackage{cprotect} % Allow \cprotect for display code with \verb|code|
\usepackage{graphicx}
\usepackage{verbatim}
\usetheme{Frankfurt}
\usecolortheme{dolphin}
\usefonttheme{professionalfonts}

\begin{document}
\author{Matemático Oromion}
\title{\Large Renombrando comandos en \LaTeX{} con}
\cprotect\subtitle{\verb|\renewcommand{cmd}[args]{def}|}
\logo{\includegraphics[scale=.5]{peru.jpg}}
\institute{\large Universidad del Perú}
\date{\textcolor{magenta}{\today}}
\subject{Fin del curso}
\setbeamercovered{transparent}
\setbeamertemplate{navigation symbols}{}

\begin{frame}[plain]
	\maketitle
\end{frame}

\begin{frame}{Contenidos}{Presentación final}
\tableofcontents
\end{frame}

\section{Introducción}
\cprotect\subsection{Entiendiendo el comando: \verb!\newcommand{\newbox}[args]{def}!}

\begin{frame}[fragile]
\frametitle{\thesection.\thesubsection. Entendiendo el comando \cmd{newcommand}}
\vspace*{-2cm}
\centering
\verb|\newcommand{cmd}[args]{def}|

\begin{columns}
\column{0.5\textwidth}
Se desea nombrar a un color.

\begin{itemize}
	\item El primer argumento obligatorio recibe el nombre del \cmd{comando}.
	\item 
\end{itemize}
\column{0.5\textwidth}
\begin{examples}{Comando \cmd{newcommand}}
\begin{enumerate}
\item[$\bullet$] A.
\end{enumerate}
\end{examples}

\begin{alertblock}{¡Cuidado!}
Se comete un error si el comando ya está definido.
\end{alertblock}
\end{columns}

\end{frame}

\cprotect\subsection{Aprendiendo el comando: \verb!\renewcommand{cmd}[args]{def}!}

\begin{frame}[fragile]
\frametitle{\thesection.\thesubsection. Aprendiendo el comando \cmd{renewcommand}}
\vspace*{-2cm}
\centering
\verb|\renewcommand{cmd}[args]{def}|

\begin{block}{Comando \cmd{renewcommand}}
Redefine un comando predefinido y comete un error si aún no está definido.
\end{block}
\end{frame}

\cprotect\subsection{Extra: \verb!\providecommand{cmd}[args]{def}!}
\cprotect\subsection{Usando el comando robusto: \verb!\@!}

\section{Ejemplos}

\subsection{Definiendo el estilo del modo matemático de una ecuación en derivadas parciales}
\subsection{Creando un comando de una caja}

\begin{frame}
\frametitle{\thesection.\thesubsection. Extra: \cmd{providecommand}}

\begin{alertblock}{Comando \cmd{providecommand}}
Define un nuevo comando si aún no está definido.
\end{alertblock}
\end{frame}

\section{\thesection. Agradecimientos}
\subsection{Ubuntu Colombia}
\begin{frame}
\frametitle{Agradecimientos}
 \end{frame}
\end{document}